% ============================================================
%  reMarkable Paper Pro Move 风格 · pandoc/xelatex 头文件
%  用法见 pandoc/remarkable.yaml。修改字体只改本文件。
% ============================================================

% ---------- 图片（mermaid 渲染图经 \includegraphics 内嵌，需显式加载）----------
\usepackage{graphicx}

% ---------- 配色（取自白皮书 HTML 的 light 主题）----------
\usepackage{xcolor}
\definecolor{rmpaper}{HTML}{ECE9E1}   % 纸张底色
\definecolor{rmraised}{HTML}{F6F4EC}  % 代码/卡片底色
\definecolor{rmink}{HTML}{201F1B}     % 正文
\definecolor{rminksoft}{HTML}{5C584D} % 次要文字
\definecolor{rminkfaint}{HTML}{8A8577}
\definecolor{rmaccent}{HTML}{2B6357}  % 主强调（雾松绿）
\definecolor{rmaccentink}{HTML}{1C4A41}
\definecolor{rmcaution}{HTML}{A85C22}
\definecolor{rmline}{HTML}{D7D2C8}    % 细线

% 纸张底色。想要纯白（打印省墨）就注释掉下面这一行。
\pagecolor{rmpaper}
\AtBeginDocument{\color{rmink}}

% ---------- 字体（项目字体，绝对路径引用，不依赖 fontconfig 注册）----------
\usepackage{fontspec}
\usepackage{xeCJK}
% 字体统一放在本地字体区（与任何项目解耦）。换机时把这个目录一起带上即可。
\def\rmfonts{/home/afu/.local/share/fonts/remarkable-pandoc/}
\def\lxgw{/home/afu/.local/share/fonts/remarkable-pandoc/}
\def\rmassets{/home/afu/.local/share/fonts/remarkable-pandoc/}

% 拉丁正文：KF Readerly（阅读衬线）
\setmainfont{KF Readerly}[
  Path=\rmfonts,
  UprightFont=KF_Readerly-Regular.ttf,
  BoldFont=KF_Readerly-Bold.ttf,
  ItalicFont=KF_Readerly-Italic.ttf,
  BoldItalicFont=KF_Readerly-BoldItalic.ttf ]

% 中文正文：霞鹜文楷（纸感楷体）
\setCJKmainfont{LXGW WenKai}[
  Path=\lxgw,
  UprightFont=LXGWWenKai-Regular.ttf,
  BoldFont=LXGWWenKai-Medium.ttf ]

% 等宽/代码：霞鹜文楷等宽（含中文，代码里的中文注释也能正常显示）
\setmonofont{LXGW WenKai Mono}[
  Path=\lxgw, Scale=0.94,
  UprightFont=LXGWWenKaiMono-Regular.ttf,
  BoldFont=LXGWWenKaiMono-Medium.ttf ]
\setCJKmonofont{LXGW WenKai Mono}[
  Path=\lxgw, Scale=0.94,
  UprightFont=LXGWWenKaiMono-Regular.ttf,
  BoldFont=LXGWWenKaiMono-Medium.ttf ]

% 标题字体：京華老宋（中文宋体显示）+ KF Readerly 粗（拉丁）
\setCJKfamilyfont{rmheadcjk}{KingHwaOldSong}[ Path=\rmassets, UprightFont=KingHwaOldSong.ttf ]
\newfontfamily\rmHeadLatin{KF Readerly}[ Path=\rmfonts,
  UprightFont=KF_Readerly-Bold.ttf, BoldFont=KF_Readerly-Bold.ttf ]
\newcommand{\rmHeadFont}{\rmHeadLatin\CJKfamily{rmheadcjk}}

% ---------- 符号/emoji 兜底（正文字体不含这些字形，映射到 Noto CJK 单色等价）----------
% MD 源里保留 emoji（网页查看好看）；PDF 用主题色单色符号呈现。
\usepackage{newunicodechar}
\newfontfamily\rmfallcjk{Noto Sans CJK SC}
\newfontfamily\rmfallsym{Noto Sans Symbols 2}
\newcommand{\rmglyph}[1]{{\rmfallcjk\char"#1}}
\newcommand{\rmsymb}[1]{{\rmfallsym\char"#1}}
\newunicodechar{→}{\rmglyph{2192}}
\newunicodechar{①}{\rmglyph{2460}}
\newunicodechar{②}{\rmglyph{2461}}
\newunicodechar{③}{\rmglyph{2462}}
\newunicodechar{④}{\rmglyph{2463}}
\newunicodechar{▶}{\rmglyph{25B6}}
\newunicodechar{▤}{\rmglyph{25A4}}
\newunicodechar{▾}{\rmsymb{25BE}}
\newunicodechar{▴}{\rmsymb{25B4}}
\newunicodechar{◑}{\rmglyph{25D1}}
\newunicodechar{✓}{\rmglyph{2713}}
\newunicodechar{✗}{\rmsymb{2717}}
\newunicodechar{⊂}{\rmglyph{2282}}
\newunicodechar{∩}{\rmglyph{2229}}
% 三个状态 emoji → 主题色单色徽标
\newunicodechar{✅}{{\color{rmaccent}\rmglyph{2713}}}
\newunicodechar{⚠}{{\color{rmcaution}\rmglyph{25B2}}}
\newunicodechar{📌}{{\color{rminksoft}\rmglyph{25C6}}}
\newunicodechar{️}{}  % U+FE0F 变体选择符，去掉

% ---------- 版式 ----------
\usepackage{geometry}
\geometry{a4paper, margin=2.2cm}   % 想更贴近设备阅读可改 a5paper
\usepackage{setspace}
\setstretch{1.32}                  % 中文行距更舒展
\usepackage{parskip}
\setlength{\parskip}{0.55em}

% ---------- 标题样式 ----------
\usepackage{titlesec}
\titleformat{\section}{\rmHeadFont\color{rmaccent}\LARGE}{}{0pt}{}
\titleformat{\subsection}{\rmHeadFont\color{rmaccentink}\Large}{}{0pt}{}
\titleformat{\subsubsection}{\bfseries\color{rmink}\large}{}{0pt}{}
\titlespacing*{\section}{0pt}{1.7em}{0.7em}
\titlespacing*{\subsection}{0pt}{1.2em}{0.5em}
\titlespacing*{\subsubsection}{0pt}{1.0em}{0.4em}
% 章节之间加一条细分隔线
\titleformat{\section}[hang]{\rmHeadFont\color{rmaccent}\LARGE}{}{0pt}
  {}[{\color{rmline}\titlerule[0.6pt]}]

% ---------- 修复溢出（核心）----------
% 1) 代码块长行折行
\usepackage{fvextra}
\DefineVerbatimEnvironment{Highlighting}{Verbatim}{%
  breaklines, breakanywhere, breaksymbolleft={}, commandchars=\\\{\}, fontsize=\small}
\fvset{breaklines, breakanywhere, fontsize=\small}
% 2) 长内联代码/路径/URL 断行（配合 filters.lua 的 \allowbreak 插入）
\usepackage{xurl}
\Urlmuskip=0mu plus 1mu\relax
\setlength{\emergencystretch}{3em}
% 3) 宽表格：缩小字号，pandoc 自身按比例分配的列宽即可自动换行
\usepackage{etoolbox}
\AtBeginEnvironment{longtable}{\small\setlength{\tabcolsep}{4pt}}
\AtBeginEnvironment{tabular}{\small}

% ---------- 代码块 / 引用块 卡片化 ----------
\usepackage{mdframed}
% 代码块：纸色底 + 细边（pandoc 定义了 Shaded 才重定义）
\ifdefined\Shaded
  \renewenvironment{Shaded}{%
    \begin{mdframed}[backgroundcolor=rmraised, linecolor=rmline, linewidth=0.4pt,
      roundcorner=3pt, innerleftmargin=8pt, innerrightmargin=8pt,
      innertopmargin=6pt, innerbottommargin=6pt, skipabove=0.7em, skipbelow=0.7em]}%
    {\end{mdframed}}
\fi
% 引用块（源码/说明卡片）：左侧强调竖条 + 纸色底
\renewenvironment{quote}{%
  \begin{mdframed}[topline=false, bottomline=false, rightline=false, leftline=true,
    linecolor=rmaccent, linewidth=2.4pt, backgroundcolor=rmraised,
    innerleftmargin=12pt, innerrightmargin=10pt, innertopmargin=7pt, innerbottommargin=7pt,
    skipabove=0.7em, skipbelow=0.7em]\small\color{rminksoft}}%
  {\end{mdframed}}

% ---------- 链接配色（延到 hyperref 加载之后）----------
\AtBeginDocument{\hypersetup{colorlinks=true, linkcolor=rmaccentink, urlcolor=rmaccent,
  citecolor=rmaccent, filecolor=rmaccent}}
